% Drop-in replacement for the vague optimisation description in Section IV.
% IMPORTANT: the public implementation is surrogate-assisted optimisation,
% not deep/offline reinforcement learning. Do not call it RL or DRL unless the
% implementation is replaced by an actual sequential decision-making agent.

\subsection{Random-Forest-Assisted Offline Optimisation}
\label{sec:rf-optimisation}

The simulator is treated as an expensive black-box function. One ns-3 run
evaluates a complete mission configuration
\begin{equation}
 \mathbf{u}=[\boldsymbol{\tau},\boldsymbol{\alpha},
 \mathbf{v}^{\min},\mathbf{v}^{\max},
 \mathbf{z}^{\min},\mathbf{z}^{\max},\boldsymbol{\Delta t}]
 \in\mathcal{U}\subset\mathbb{R}^{28},
 \label{eq:decision-vector}
\end{equation}
where the four entries in each group correspond to the surveillance, patrol,
rapid-response, and strategic UAV roles. The groups respectively contain the
sensing intervals, Gauss--Markov memory coefficients, minimum and maximum
speeds, minimum and maximum altitudes, and mobility-update intervals. The run
returns the detection probability $P_{\rm det}$, throughput $T$, localisation
error $R$, delay $D$, packet-loss percentage $L$, and energy consumption $E$.

To combine quantities having different units, fixed operational reference
values are used rather than adding the raw metrics. With
$c(x)=\min(1,\max(0,x))$, the implementation computes
\begin{equation}
 J(\mathbf{u})=3c(P_{\rm det})+3c(T/0.002)
 -2c(R/30)-c(D/20)-c(L/100),
 \label{eq:implemented-objective}
\end{equation}
where $T$, $R$, and $D$ are expressed in Mbit/s, metres, and milliseconds,
respectively. Thus every term is dimensionless and bounded before weighting.
The weights prioritise successful detection and delivery of sensing data,
penalise localisation error more strongly than latency or loss, and remain
fixed in every experiment. The released optimiser reports energy separately
but does not use it when ranking candidates. Hence, in the reported
implementation, battery feasibility is an a-posteriori check rather than an
enforced search constraint; a hard constraint claim requires candidates with
$E_i(\mathbf{u})>B_i$ to be rejected before selection.

After $K_0=40$ space-filling baseline runs, the observations form
$\mathcal{D}_0=\{(\mathbf{u}_k,J_k)\}_{k=1}^{K_0}$. A random-forest surrogate
$\widehat{J}(\mathbf{u})$ is fitted using 300 trees, maximum tree depth 8,
minimum leaf size 2, and random seed 42. Model adequacy is checked using
training MAE and $R^2$ together with leave-one-out cross-validation MAE. At
optimisation round $r$, 50,000 configurations are sampled uniformly inside
the prescribed parameter bounds. Infeasible structural combinations are
repaired to enforce $v_i^{\max}\geq v_i^{\min}+1$~m/s and
$z_i^{\max}\geq z_i^{\min}+5$~m; previously evaluated and duplicate
configurations are removed. The surrogate ranks the remaining candidates,
and the five candidates having the largest predicted score are evaluated in
ns-3. Their measured scores, rather than their predictions, are appended to
the dataset before the surrogate is refitted. The process terminates when the
simulation budget is exhausted or the best measured score ceases to improve.
Algorithm~\ref{alg:rf-search} summarises this procedure.

\begin{algorithm}[t]
\caption{Random-forest-assisted UAV configuration search}
\label{alg:rf-search}
\begin{algorithmic}[1]
\REQUIRE Bounds $\mathcal{U}$, initial budget $K_0=40$, candidate pool
$N_c=50{,}000$, batch size $q=5$
\STATE Evaluate $K_0$ baseline configurations in ns-3 to obtain
$\mathcal{D}_0$
\FOR{$r=0,1,\ldots$ until the stopping condition is met}
  \STATE Fit the 300-tree random-forest surrogate $\widehat{J}_r$ to
  $\mathcal{D}_r$
  \STATE Sample $N_c$ candidates in $\mathcal{U}$ and enforce the implemented
  speed and altitude constraints
  \STATE Remove duplicate and previously evaluated candidates
  \STATE Select the $q$ candidates with the largest
  $\widehat{J}_r(\mathbf{u})$
  \STATE Evaluate the selected candidates in ns-3
  \STATE Append their measured $(\mathbf{u},J)$ pairs to obtain
  $\mathcal{D}_{r+1}$
\ENDFOR
\RETURN $\mathbf{u}^{\star}=\arg\max_{(\mathbf{u},J)\in\mathcal{D}}J$
\end{algorithmic}
\end{algorithm}

This method is used instead of model-free reinforcement learning because a
simulation run produces one terminal score for a complete mission
configuration rather than a sequence of state--action transitions. It also
avoids requiring gradients and is suitable for a small tabular dataset with
nonlinear parameter interactions. Accordingly, the method is more precisely
classified as surrogate-assisted black-box optimisation, not as deep or
offline reinforcement learning; an MDP formulation, policy network, discount
factor, and Bellman update do not apply to the implemented method.

\subsection{Equal-Budget Benchmark Protocol}
\label{sec:optimisation-benchmarks}

To separate gains due to the surrogate from gains due merely to additional
simulation trials, the proposed search must be compared with independent
uniform random search and at least one established black-box optimiser (e.g.,
Gaussian-process Bayesian optimisation or a genetic algorithm). Every method
starts from the same $K_0$ observations, receives the same number of additional
ns-3 evaluations, and is repeated with the same set of simulator seeds. The
comparison reports the best measured feasible score versus evaluation count,
final $P_{\rm det}$, RMSE, throughput, delay, loss, and energy, together with
mean and 95\% confidence intervals across repetitions. This equal-budget
design avoids comparing 40 unguided trials with 14 adaptively selected trials
solely through phase averages.

% Insert measured results only after running the above protocol:
% ``After [N] independent repetitions, RF-assisted search achieved [...],
% compared with [...] for random search and [...] for [BO/GA].''
